% INTRODUCTION. Numbers ONLY via \R... macros from generated/numbers.tex.
An LLM agent's prompt is assembled, not written. Before each decision a harness collects tool schemas, retrieved
documents, worked examples, notes and conversation history, and must fit them into a token budget set by latency, cost or
the model's window~\cite{cesurvey2025}. What to keep is usually decided by heuristics: a relevance score against the
goal, a coverage objective over candidate passages~\cite{pacms2026}, a compressor that deletes tokens~\cite{llmlingua2_2024},
or fixed priorities written by an engineer. None of these asks the question a budgeted system actually faces:
\emph{how much does this kind of context change task success for this model on this task family?} The answer is not
obvious. Models use long contexts unevenly~\cite{lostmiddle2024}, and a block can be relevant to the goal yet add
nothing the model needs, or be redundant with another block that is already included.

This paper asks: when context consists of heterogeneous types, can we measure how much each type contributes to task
success, and can those measurements be reused for budgeted assembly? We separate three levels of the problem
(Section~\ref{sec:problem}): \emph{type value} (which kinds of context matter for a model and task family),
\emph{block relevance} (which particular blocks matter for an instance), and \emph{budget allocation} (what to send under a
token limit, and when to stop). A \emph{marginal context value (MCV) profile} answers the first level. For a (model, task
family) pair it records the measured drop in success when every block of a type is removed, the value of the type's
short variant, and pairwise interactions between types, each with a bootstrap interval. Profiles are estimated once,
offline, by counterfactual ablation on development instances, with a length-neutral control. For the second level we
train a block-value model on leave-one-block-out deltas from the same runs. For the third, a CP-SAT compiler maximizes
predicted value under the budget, dependencies and variants without calling the language model, and reports a
certificate of optimality for that predicted objective.

Prior work values context in related ways but for other decisions: attribution methods explain one generation by
ablating its sources~\cite{contextcite2024,attribot2024}, Shapley-style methods value demonstrations, prompts, data
points or skills~\cite{demoshapley2024,promptshapley2023,datashapley2019,skillsv2026}, and recent context-allocation
studies measure the causal effect of context in generative search~\cite{ctxalloc2026}. Section~\ref{sec:related}
positions MCV profiles against these and against budgeted selectors. The distinguishing combination is a reusable,
type-level, interval-reported value table measured per model and task family and compiled under a hard token budget.

We evaluate the approach under a pre-registered protocol (hypotheses, margins, sample sizes and the method version fixed
before any test run, recorded as a commit in the released repository; it was not deposited with an external registry) on three task families: multi-hop question answering over HotpotQA documents with distractors,
tool selection and argument filling over a \RNumTools-tool catalog built from BFCL, and recall of the latest value of a
user attribute from \RNumSessions synthetic sessions. All models run locally (Qwen3-4B-Instruct for the main study,
Qwen3-8B for transfer). We compare against six baselines, including LLMLingua-2 and a PACMS-style coverage selector.

The results are mixed, and we report them as such. \textbf{The primary hypothesis failed:} MCV-compiled context did not
achieve a higher area under the success--budget curve than the best baseline (embedding relevance) in any family. It was
\RHOneHistoryAbs{} lower on history, where at the smallest budget the block-value model kept the session containing the
answer in \RRecallOurs of plans, against \RRecallRel for relevance ranking. \textbf{The efficiency hypothesis}, added after the development pilots but before any test run, \textbf{held} in
\RHOnebPass of three families. At an 8k budget the compiler, which leaves budget unused when the remaining blocks add no
predicted value, used \RTokCutTool fewer prompt tokens than full context on tool use and \RTokCutHistory fewer on history
with non-inferior success. On question answering it used \RTokCutFacts fewer, but non-inferiority was not shown.
Interaction terms had no measurable effect. The profiles themselves were informative. Tool schemas and documents carry
large value (\RMcvToolToolSchema and \RMcvFactsFact), while worked examples carry little value on their own in all three families.
History and the profile note substitute for each other (interaction \RIntHistoryHistoryTurnNote). The length-neutral control moved no
value by more than \RPadDiffMax, and the type values of the 4B and 8B models were rank-correlated ($\rho=\RRho$).

\textbf{Contributions.}
\begin{itemize}
\item \emph{A pre-registered, fully local evaluation} of eight context policies on three task families with measured
token counts, paired tests and Holm correction. Its main findings are negative or qualified: measured context value did
not beat relevance ranking for selection, and single-block deletions proved too sparse a signal to train a block
ranker that improves on embedding similarity (Section~\ref{sec:eval}).
\item \emph{MCV profiles as a diagnostic:} a definition and estimator of per-(model, family) type values,
short-variant values and pairwise interactions with shrinkage, bootstrap intervals and a length-neutral control,
which identified low-value context types and, in subsampling, stabilized with ten to twenty development instances
(Sections~\ref{sec:problem}, \ref{sec:design} and~\ref{sec:eval}).
\item \emph{A budget compiler with certificates:} a CP-SAT program over blocks, variants, dependencies and type
interactions that enforces the budget in the model's own tokenizer and reports the solver's status, objective and bound.
The certificate covers optimality with respect to the predicted values only, not task success
(Section~\ref{sec:design}).
\item \emph{An open artifact:} the \texttt{mcv} package, data adapters, logs of every model call, and the scripts that
produce every number and figure in this paper from those logs.
\end{itemize}
